% ------------------------------------------------------------
% CHAPTER 23
% ------------------------------------------------------------

\chapter{Delta–Sigma as Active Inference}

The connection between delta–sigma modulation and active inference becomes evident once we observe that both systems are organized around the minimization of prediction error under constraints imposed by discrete symbolic collapse. In active inference, an agent maintains internal beliefs about the external world and continually updates these beliefs by minimizing a free-energy functional that balances accuracy, complexity, and uncertainty. In delta–sigma modulation, the internal state of the loop is continually updated to minimize the discrepancy between the physical input and its quantized representation, with quantization acting as the source of uncertainty. The structural correspondence is so exact that the delta–sigma loop may be viewed as the simplest nontrivial active inference engine. Its dynamics instantiate the universal architecture of prediction, collapse, correction, and coherence preservation.

To formalize this relationship, let \(u(t)\) denote the external world state, represented as a continuous scalar field. The delta–sigma loop maintains an internal state \(x(t)\), which plays the role of a belief about the world. The quantizer output \(y[n]\) corresponds to a symbolic percept, akin to a discrete observation in Bayesian inference. In active inference, agents generate predictions of sensory input, compare them to actual input, compute prediction errors, and use these errors to update internal states. Similarly, in the delta–sigma modulator, the loop filter computes a prediction \(x[n]\) of the quantizer input, and the discrepancy between \(u[n]\) and \(Q(x[n])\) yields the error signal that drives the loop.

This correspondence becomes precise when we introduce the variational free-energy functional
\[
\mathcal{F}(x) = \frac{1}{2} \int \left( u(t) - x(t) \right)^{2} dt + \lambda \int S(t)\, dt,
\]
which measures the discrepancy between the true field \(u\) and its internal representation \(x\), penalized by the entropy field \(S\) associated with quantization. Minimizing \(\mathcal{F}\) is equivalent to performing inference under uncertainty. In active inference, the entropy term reflects uncertainty in the agent’s generative model. In delta–sigma modulation, the entropy term arises from quantization noise, which is injected by the collapse operator \(Q\). The delta–sigma loop performs gradient descent on this functional by updating the internal state according to a relation of the form
\[
x[n+1] = x[n] - \eta\left( x[n] - u[n] + \lambda \frac{\partial S}{\partial x} \right),
\]
where the derivative of the entropy term captures the quantization-induced perturbation. This equation is structurally identical to the update rule governing predictive coding networks and free-energy minimization engines.

An essential feature of active inference is the existence of a generative model, which predicts the sensory input based on internal dynamics. In delta–sigma modulation, the loop filter plays the role of the generative model. For a first-order modulator, the generative model assumes that the input signal is slowly varying, which is implemented by an integrator. Higher-order modulators assume that higher-order derivatives of the input vanish or remain bounded. This assumption defines the curvature of the belief manifold. The vector field associated with the loop filter therefore defines a prior over the trajectories of the field. If the loop filter has transfer function \(H(z)\), then the prior can be interpreted as the assumption that the field obeys the differential equation
\[
H\left(\frac{d}{dt}\right)\Phi(t) = 0.
\]
The modulator attempts to infer the field \(\Phi\) that satisfies this prior while remaining consistent with the quantized data. This interplay between prediction and correction is the essence of inference.

The quantizer acts as a symbolic perception mechanism. In Bayesian inference, the likelihood function maps hidden states to observations with associated uncertainty. In delta–sigma modulation, the quantizer maps internal states to symbols by performing a discontinuous projection onto a discrete set. The discrepancy between the predicted symbol and the observed symbol generates a prediction error, which flows backward through the vector field and corrects the internal state. In this sense, the quantizer output \(y[n]\) is analogous to perceptual evidence in cognitive systems. The perception mechanism is lossy and quantized, and the feedback loop performs inference to reconstruct the underlying continuous causes.

A key insight emerges when we analyze the error dynamics. In delta–sigma modulation, the quantization error is shaped by the noise transfer function, which ensures that most of the error energy is routed into high-frequency modes. In active inference, prediction errors are similarly decomposed into components that are either used to update beliefs or disregarded because they are inconsistent with the generative model. The delta–sigma loop implements an implicit decomposition of error into modes that are relevant to the belief update and modes that are irrelevant. High-frequency errors are deemed irrelevant because they violate the assumption of low curvature; thus they are suppressed. The modulator therefore performs a form of epistemic filtering, attributing certain discrepancies to noise rather than structure.

This understanding becomes even clearer when we examine the continuous-time limit. In the limit of infinite oversampling ratio, the delta–sigma modulator approaches a differential equation of the form
\[
\frac{d^{L} x}{dt^{L}} = u(t) - Q(x(t)),
\]
which is a gradient flow on the free-energy functional. The right-hand side represents prediction error. When the quantizer collapses the state, it introduces a discontinuity, which the system smooths through lamphrodynamic relaxation. Thus the delta–sigma loop continually seeks a geodesic on the belief manifold while compensating for discrete disturbances. This is identical to the dynamics of active inference, in which the system follows a geodesic flow on an information manifold perturbed by observations.

The connection becomes even more striking when we examine the role of stabilizing potentials. In active inference, variational free energy decomposes into accuracy and complexity terms. The complexity term penalizes deviations of beliefs from the prior. In delta–sigma modulation, stability is ensured by shaping the internal dynamics so that the state remains within a bounded region of the field configuration space. This is achieved through implicit potentials arising from the loop filter. In high-order modulators, one must design these potentials carefully to avoid runaway behavior. This corresponds to the complexity term in active inference, which prevents the model from adopting beliefs that are too complex relative to the available data.

Finally, we observe that delta–sigma modulators can be interpreted as minimal agents. They possess internal states that encode beliefs, they generate predictions, they receive observations, and they adjust their beliefs to minimize prediction error while redistributing entropy. The dynamics of the loop embody agency at the simplest level: the maintenance of coherence despite discrete collapse. This perspective suggests that delta–sigma modulators can serve as building blocks for artificial agents that combine sensing, inference, and action within a unified feedback architecture.

Thus, delta–sigma modulation emerges not merely as a circuit design technique but as a concrete realization of the principles governing active inference. The modulator’s structure implements prediction, collapse, and correction through dynamical processes that minimize a free-energy functional, shape entropy, and maintain coherence across time. This unification prepares the conceptual ground for the next chapter, which examines quantization through the lens of semantic collapse and symbolic discretization.


